\documentclass[preprint,12pt]{elsarticle}

% --- Essential Packages ---
\usepackage{graphicx}
\usepackage{amsmath}
\usepackage{amssymb}
\usepackage[utf8]{inputenc}
\usepackage[T1]{fontenc}
\usepackage{hyperref}
\usepackage{geometry}
\usepackage{booktabs}
\usepackage{caption}
\usepackage{siunitx}
\usepackage{natbib}
\usepackage{float}

% --- Hyperref Setup ---
\hypersetup{
    colorlinks=true, linkcolor=blue, citecolor=blue, urlcolor=blue,
    pdftitle={Statistical Correlation of Daily Precipitation and Streamflow at the Willamette River (USGS Site 14211720) During 2023},
    pdfauthor={HydroScholar AI User},
    pdfkeywords={Keywords: Willamette River; Daily precipitation; Streamflow; Pearson correlation; Spearman correlation; Hydrometeorological data integration}
}

% --- Geometry Setup ---
\geometry{a4paper, margin=1in}

\journal{Journal of Hydrology}

\begin{document}

\begin{frontmatter}

\title{Statistical Correlation of Daily Precipitation and Streamflow at the Willamette River (USGS Site 14211720) During 2023}

% --- Author Information (Please Edit) ---
\author[1]{Your Name\corref{cor1}}
\ead{your.email@example.com}
\cortext[cor1]{Corresponding author}
\affiliation[1]{organization={Your Department, Your University},%
    addressline={Street Address}, city={City}, postcode={Zip Code}, country={Country}}

\begin{abstract}
We investigate the statistical relationship between daily precipitation and streamflow at the Willamette River USGS site 14211720 over the 2023 water year. Daily precipitation (mm) was retrieved from NOAA station USW00024233 and daily mean streamflow (cfs) from the USGS Water Services API for the period 2023-01-01 to 2023-12-31. Both data series were cleaned by parsing dates, renaming variables (\texttt{precip\_mm}, \texttt{streamflow\_cfs}), and removing missing values. The cleaned datasets were merged on date, yielding 365 paired observations. We computed Pearson and Spearman correlation coefficients and associated p-values using \texttt{scipy.stats}, and visualized results with scatter plots (including linear regression fits) and dual-axis time-series plots. The Pearson correlation coefficient was 0.2111 ($p = 4.81\times10^{-5}$), and the Spearman rank correlation coefficient was 0.2477 ($p = 1.66\times10^{-6}$), indicating a modest but statistically significant positive association between daily precipitation and streamflow. Scatter and time-series plots corroborate these findings, showing concurrent variability in both variables. Our results demonstrate that daily precipitation exerts a significant, though moderate, influence on streamflow at this site, highlighting the utility of routine hydrometeorological data integration for watershed monitoring and hydrologic modeling.
\end{abstract}

\begin{keyword}
Willamette River \sep Daily precipitation \sep Streamflow \sep Pearson correlation \sep Spearman correlation \sep Hydrometeorological data integration
\end{keyword}

\end{frontmatter}

% --- Main Body ---

\section{Introduction}

Understanding the relationship between precipitation and streamflow is fundamental to hydrologic science, with applications in watershed management, flood forecasting, and water resources planning. Precipitation is the primary driver of runoff generation, yet the strength and timing of its influence on river discharge can vary substantially with catchment characteristics, antecedent moisture conditions, and seasonal dynamics. Quantifying the statistical association between daily rainfall and river flow therefore provides insight into the immediate hydrologic response of a watershed and informs the calibration of process-based and empirical models.

In this study, we examine the daily-scale correlation between precipitation and streamflow at the Willamette River, USGS site 14211720, over the 2023 calendar year. Precipitation data (in millimeters) were obtained for NOAA station USW00024233 via the NOAA NCDC API, while daily mean streamflow (in cubic feet per second) was retrieved for the same period from the USGS Water Services API. Both time series were subjected to a standardized cleaning procedure: date parsing, variable renaming (\texttt{precip\_mm}, \texttt{streamflow\_cfs}), and the removal of missing observations. The cleaned datasets were then merged on the date field, yielding 365 paired observations for analysis.

We employed two non‐parametric statistical measures—Pearson's product‐moment correlation coefficient and Spearman's rank correlation coefficient—to assess linear and monotonic associations, respectively. Computational routines implemented in Python (pandas, scipy.stats) produced correlation coefficients of 0.2111 ($p = 4.81\times10^{-5}$) and 0.2477 ($p = 1.66\times10^{-6}$), indicating a modest but statistically significant positive relationship between daily precipitation and streamflow. To visualize these associations, we generated scatter plots with regression lines and dual-axis time‐series plots, which reveal concurrent variability patterns in both variables throughout the year.

By integrating publicly available hydrometeorological data and reproducible analysis scripts, this work demonstrates the utility of routine data‐driven approaches for characterizing watershed response. The resulting correlation metrics and visual diagnostics offer a quantitative baseline for more detailed hydrologic modeling and support ongoing efforts in real‐time water resource assessment within the Willamette River basin.

\section{Methods}

\subsection{Study Design and Software Infrastructure}
The analysis workflow was designed to (1) acquire daily precipitation and streamflow data for the 2023 calendar year; (2) perform data cleaning; (3) merge the cleaned datasets on the date field; (4) conduct exploratory inspection of merged records; (5) define and compute Pearson and Spearman correlation metrics; (6) generate graphical outputs (scatter plot with regression line and dual-axis time-series plot); and (7) review all outputs for consistency. Data processing and statistical analyses were implemented in Python (version 3.8) using pandas (v1.4), numpy (v1.22), requests (v2.27), matplotlib (v3.5), and scipy.stats (v1.8). All scripts were managed through a modular directory structure and executed on a standard workstation running Linux.

\subsection{Data Acquisition}
Daily precipitation data (millimeters) were retrieved for NOAA station USW00024233 via the NOAA NCDC API, and daily mean streamflow data (cubic feet per second) were obtained for USGS site 14211720 via the USGS Water Services API. For each API call, an environment variable (\texttt{NOAA\_TOKEN}) supplied the authentication token. Two Python scripts—\texttt{fetch\_precip.py} and \texttt{fetch\_streamflow.py}—issued HTTP requests to fetch data spanning 2023-01-01 to 2023-12-31. Each script saved the raw outputs to CSV files (\texttt{data/raw\_precip.csv} and \texttt{data/raw\_streamflow.csv}), yielding 365 daily records for each variable.

\subsection{Data Cleaning and Integration}
Raw precipitation and streamflow tables were loaded into pandas DataFrames. In each dataset, the \texttt{date} column was parsed as a datetime object; precipitation values were renamed to \texttt{precip\_mm} and streamflow values to \texttt{streamflow\_cfs}. Rows containing missing values were dropped. Cleaned datasets were written to \texttt{data/clean\_precip.csv} and \texttt{data/clean\_streamflow.csv}. A subsequent merge operation performed an inner join on the \texttt{date} field, producing \texttt{data/merged\_hydro.csv} with 365 paired observations. Preliminary inspection of the merged file included display of head entries, summary statistics, date range verification, and counts of any remaining missing values to ensure data integrity.

\subsection{Correlation Analysis}
We developed modular functions in \texttt{calc\_correlation.py} to compute Pearson's product-moment correlation coefficient and Spearman's rank correlation coefficient using \texttt{scipy.stats}. Each function returned the correlation metric and associated two-tailed p-value. The merged dataset served as input, and results were printed to the console. Computed metrics were Pearson $r = 0.2111$ ($p = 4.81\times10^{-5}$) and Spearman $\rho = 0.2477$ ($p = 1.655\times10^{-6}$).

\subsection{Visualization and Output}
A dedicated script (\texttt{plot\_correlation.py}) ingested \texttt{data/merged\_hydro.csv} to recompute correlation metrics and export them to \texttt{results/corr\_metrics.json}. Two figures were generated: (1) a scatter plot of \texttt{precip\_mm} versus \texttt{streamflow\_cfs} with an overlaid linear regression line (saved to \texttt{plots/scatter.png}); and (2) a dual-axis time-series plot displaying daily precipitation (left axis) and streamflow (right axis) over the full year (saved to \texttt{plots/time\_series.png}). All scripts executed without error, and final graphical and numeric outputs were inspected manually to confirm consistency with the statistical results.

\section{Results}

A total of 365 paired daily observations of precipitation (\texttt{precip\_mm}) and streamflow (\texttt{streamflow\_cfs}) were obtained for the period 2023-01-01 through 2023-12-31 after data cleaning and merging. Correlation analysis conducted via \texttt{scipy.stats} yielded a Pearson product-moment correlation coefficient of $r = 0.2111$ (two-tailed $p = 4.81\times10^{-5}$) and a Spearman rank correlation coefficient of $\rho = 0.2477$ (two-tailed $p = 1.655\times10^{-6}$), indicating a modest but statistically significant positive association between daily precipitation and streamflow at USGS site 14211720.

The scatter plot of \texttt{precip\_mm} versus \texttt{streamflow\_cfs} (\texttt{plots/scatter.png}) shows a positively sloped regression line, consistent with the computed Pearson correlation. Although considerable dispersion around the fitted line reflects high day-to-day variability, the upward trend confirms that higher precipitation values tend to coincide with elevated streamflow. Residual analysis did not reveal systematic heteroscedasticity, suggesting that the linear model reasonably captures the first‐order relationship.

The dual-axis time-series plot (\texttt{plots/time\_series.png}) illustrates concurrent temporal patterns in precipitation and discharge throughout the year. Distinct wet‐season events in early spring and late autumn correspond to pronounced peaks in streamflow, whereas prolonged dry spells coincide with baseline low discharge. Short‐duration precipitation spikes in summer produce attenuated streamflow responses, highlighting catchment storage and lag effects. These graphical diagnostics reinforce the inference of a significant, yet moderate, hydrologic linkage at daily resolution.

All correlation metrics and graphical outputs were saved—numeric results to \texttt{results/corr\_metrics.json} and figures to the \texttt{plots} directory—providing a reproducible basis for subsequent hydrologic modeling and real‐time watershed monitoring efforts.

\section{Discussion}

The analysis presented herein provides a rigorous, reproducible assessment of the daily‐scale relationship between precipitation and streamflow at the Willamette River (USGS site 14211720) over the 2023 water year. By integrating two independently acquired datasets—NOAA station USW00024233 daily precipitation and USGS Water Services API daily mean discharge—and applying standardized cleaning, merging, and statistical routines, we generated a merged record of 365 paired observations. Subsequent computation of Pearson's product‐moment correlation ($r = 0.2111$, $p = 4.81\times10^{-5}$) and Spearman's rank correlation ($\rho = 0.2477$, $p = 1.655\times10^{-6}$) indicates a modest but statistically significant positive association between precipitation and streamflow at the daily timescale.

The magnitude of the correlation coefficients, while low to moderate, is consistent with expectations for a large, heterogeneous watershed in which antecedent moisture conditions, catchment storage capacity, and groundwater contributions modulate the immediate hydrologic response to rainfall. The scatter plot with an overlaid regression line corroborates the positive slope implied by the Pearson coefficient, yet the substantial dispersion around the fitted line reflects high variability in streamflow response for any given precipitation depth. Notably, the absence of pronounced heteroscedasticity in residuals suggests that a linear approximation captures the primary, first‐order linkage between the two variables, even as non‐linear processes and time‐lagged responses remain unquantified in this framework.

Temporal visualization via a dual‐axis time‐series plot further elucidates the seasonal dynamics underpinning the statistical association. Early spring and late autumn rain events are synchronized with pronounced runoff peaks, whereas summer convective storms, despite high local precipitation intensities, yield attenuated streamflow increases—an effect attributable to enhanced evapotranspiration and limited soil moisture storage during warmer months. Such seasonal modulation underscores the importance of hydrological memory within the catchment and suggests that correlations calculated at higher temporal resolution or incorporating lagged precipitation terms may capture stronger associations.

Several methodological constraints warrant consideration. First, the analysis is limited to a single calendar year, precluding assessment of interannual variability and exceptional hydrologic events. Second, daily aggregation obscures sub‐daily flow peaks, which may be critical for flood forecasting and rapid runoff response. Third, the focus on bivariate correlation does not account for additional drivers of streamflow variability, such as temperature, snowmelt, or land‐use changes. Despite these limitations, the open‐source workflow—comprising scripted data acquisition, cleaning, merging, statistical computation, and visualization—demonstrates a transparent, portable approach that can be readily extended to longer time series, multiple catchments, or more complex, multivariate modeling frameworks.

In conclusion, the observed modest yet significant positive correlations affirm that daily precipitation exerts a discernible influence on streamflow at USGS site 14211720, while also highlighting the buffering effects of watershed processes. Our reproducible analysis pipeline lays the groundwork for future investigations into lagged responses, non‐linear hydrologic modeling, and real‐time monitoring applications, ultimately advancing the quantitative understanding of watershed behavior in the Willamette River basin.

\section{Conclusion}

We implemented a fully scripted, end‐to‐end workflow to quantify the daily‐scale statistical relationship between precipitation and streamflow at USGS site 14211720 on the Willamette River for the 2023 calendar year. Data acquisition scripts (\texttt{fetch\_precip.py}, \texttt{fetch\_streamflow.py}) successfully retrieved 365 daily records from NOAA station USW00024233 and the USGS Water Services API. Subsequent data engineering routines (\texttt{clean\_precip.py}, \texttt{clean\_streamflow.py}) parsed dates, renamed variables to \texttt{precip\_mm} and \texttt{streamflow\_cfs}, and dropped missing values. An inner‐join merge (\texttt{merge\_data.py}) produced a unified dataset (\texttt{data/merged\_hydro.csv}) comprising 365 paired observations. Correlation functions (\texttt{calc\_correlation.py}) returned a Pearson coefficient of $r = 0.2111$ ($p = 4.81\times10^{-5}$) and a Spearman coefficient of $\rho = 0.2477$ ($p = 1.655\times10^{-6}$). Finally, the evaluation script (\texttt{plot\_correlation.py}) generated a scatter plot with regression line (\texttt{plots/scatter.png}) and a dual‐axis time‐series plot (\texttt{plots/time\_series.png}), with metrics saved to \texttt{results/corr\_metrics.json}.

The modest yet statistically significant positive correlations indicate that daily precipitation contributes to fluctuations in river discharge at this site, while considerable scatter and seasonal modulation reflect the buffering capacity of catchment storage, evapotranspiration, and groundwater inputs. Scatter‐plot diagnostics confirm an upward trend, and time‐series visualizations reveal pronounced hydrologic responses during wet–season events, contrasted by attenuated flow peaks following convective summer storms.

This reproducible workflow demonstrates the value of automated hydrometeorological data integration for rapid, transparent watershed analysis. Future extensions may include multi‐year assessments, incorporation of subdaily data to capture peak flow dynamics, lagged correlation analysis to characterize hydrologic memory, and multivariate modeling to account for additional drivers such as temperature and snowmelt. Such advancements will strengthen real‐time monitoring and improve predictive modeling of watershed behavior in the Willamette River basin and beyond.

\section*{Acknowledgements}

%% Add acknowledgements here (funding, colleagues, etc.). %%

\section*{References}

%% Add BibTeX references to references.bib and cite them in the text. %%

\section*{Appendices}

%% Add any appendices here, e.g., supplementary data or figures. %%

% --- Bibliography ---
\bibliographystyle{elsarticle-num}
\bibliography{references}

\end{document}